\chapter{Conclusion}
\label{chap:discussion}

The present work provides a unified picture of ultrafast electronic dynamics in disordered silicon, combining real-time TDDFT simulations with analysis of expressions from the semiconductor Bloch equations and Drude transport theory. By systematically varying both structural disorder and excitation regime, we have identified how disorder modifies (i) the electronic structure, (ii) photoinjection efficiency, (iii) intraband transport, and (iv) coherent interband dynamics.

\section*{Electronic Structure as the Microscopic Origin}

The density-of-states analysis establishes the microscopic foundation for all subsequent dynamical effects. With increasing disorder, three key modifications occur: smoothing of van Hove singularities, systematic band-gap reduction, and formation of near-exponential band-edge tails consistent with Urbach-type behavior \cite{MottDavis1979,Urbach1953}. 

Band-gap narrowing explains the red shift observed in high-harmonic spectra and coherent current oscillation (CCO) spectra. At the same time, the emergence of tail states increases the density of available transition frequencies near the band edge, thereby enhancing inhomogeneous dephasing. Importantly, the Urbach energy extracted from the DOS correlates with the momentum relaxation time obtained from transport analysis, indicating that both spectral broadening and carrier scattering originate from the same disorder-induced potential fluctuations.

\section*{Photoinjection Regimes and Nonlinearity}

The photoinjection analysis demonstrates that excitation regime plays a central role in shaping subsequent dynamics. When the pump photon energy approaches the effective band gap, excitation enters a single-photon regime. In this limit, photoinjection becomes less sensitive to disorder-induced variations in local structure, and spectral phase distributions remain relatively narrow.

In contrast, multiphoton excitation produces broader distributions in $\mathbf{k}$-space. This leads to stronger inhomogeneous dephasing even in the absence of scattering. Disorder accelerates the transition from multiphoton-dominated excitation toward effectively linear absorption by relaxing momentum selection rules, consistent with general disorder-induced breakdown of Bloch momentum conservation \cite{Sheng2006Localization}.

\section*{Transport: Momentum Relaxation and Effective Mass Renormalization}

The probe-direction current analysis reveals that intraband transport is well described by a generalized Drude model with a disorder-controlled momentum relaxation time $\tau_m$ and a $\mathbf{k}$-displacement-dependent effective mass $m^*(k)$. 

The key observation is that $\tau_m$ depends primarily on disorder strength and only weakly on pump photon energy or probe-induced $\mathbf{k}$ displacement. The extracted relaxation times decrease monotonically with disorder, reflecting increased scattering probability due to stronger potential fluctuations.

% This indicates that elastic impurity scattering dominates the transport decay channel, in agreement with semiclassical scattering theory \cite{AshcroftMermin,Mahan2000}.

At the same time, the effective mass exhibits significant variation with reciprocal-space displacement, especially in the low-disorder regime. As disorder increases, the range of $m^*$ variation decreases. Physically, this can be understood as disorder-induced averaging over band curvature due to relaxation of strict momentum selection, similar to trends observed in strong-field excitation \cite{Qasim2018}.

% similar to trends observed in strong-field excitation studies where more uniform band population reduces effective-mass sensitivity to excitation regime \cite{Qasim2018}.

\section*{Coherent Current Oscillations: Dephasing Mechanisms}

The CCO analysis provides the clearest separation between homogeneous and inhomogeneous dephasing mechanisms. In the absence of disorder, decay is governed entirely by phase mixing across the excited $\mathbf{k}$ region, yielding an inhomogeneous dephasing time $T_{\mathrm{inh}}$. No microscopic irreversible process is present in RT-TDDFT, and therefore $T_2=\infty$ in the absence of disorder.

With static disorder, elastic scattering introduces a finite homogeneous dephasing time $T_2$. The observed envelope of coherent current oscillations is well described by a product of exponential (homogeneous) and Gaussian-like (inhomogeneous) contributions. 

A central result of this work is that the homogeneous component of CCO decay is quantitatively consistent with the transport relaxation time extracted independently from probe-direction dynamics. Within error bars, for almost all points the decay-rate additivity relation

\begin{equation}
\frac{1}{\tau_{\mathrm{coh}}^z}
=
\frac{1}{\tau_m}
+
\frac{1}{\tau_{\mathrm{coh}}^{z,\mathrm{norm}}}
\end{equation}

is satisfied. This implies that the scattering responsible for momentum relaxation also governs interband coherence decay. In terms of scattering theory, this indicates that the angular factor in the transport-time definition is close to unity, corresponding to nearly isotropic impurity scattering in the displacement-disorder model.

% \section*{Interplay of Excitation Regime and Disorder}

% An important qualitative trend emerges when comparing different pump photon energies. In the single-photon regime, coherent oscillations are spectrally cleaner and less chaotic. In the strongly nonlinear regime, the initial phase distribution is broader and more irregular, amplifying inhomogeneous contributions even before disorder is introduced.

% Disorder therefore acts on two levels. First, it modifies the band structure and transition frequencies, thereby changing the spectral content of coherent oscillations. Second, it introduces elastic scattering, which generates homogeneous exponential decay and suppresses both interband coherence and intraband current.

% The relative importance of these mechanisms depends on the excitation regime. In near-single-photon excitation, homogeneous scattering dominates the decay. In deeply nonlinear excitation, inhomogeneous phase dispersion can remain significant even when disorder is moderate.

\section*{Unified Physical Picture}

Taken together, the results support a consistent physical interpretation:

1. Structural disorder reshapes the electronic density of states, producing band-gap reduction and near-exponential band-edge tails.

2. These structural modifications determine both photoinjection efficiency and the spectral distribution of excited states.

3. Elastic scattering from the same disorder potential governs both momentum relaxation and homogeneous coherence decay.

4. Inhomogeneous dephasing arises from dispersion of transition frequencies across the excited $\mathbf{k}$ region and is enhanced by nonlinear excitation.

5. Intraband transport can be accurately described by a generalized Drude model, incorporating a disorder-induced momentum relaxation time $\tau_m$ and an effective mass $m^*(k)$ that depends on the $\mathbf{k}$-space displacement.

Thus, transport time, homogeneous dephasing time, and Urbach energy are not independent quantities but manifestations of the same underlying disorder potential.

\section*{Implications}

The demonstrated equivalence between the homogeneous component of CCO decay and the transport relaxation time suggests that ultrafast coherent-current measurements can provide direct access to transport properties on femtosecond time scales. This connects attosecond solid-state spectroscopy \cite{Schultze2014Si,Lucchini2016Science} with semiclassical transport theory in a quantitative manner.

Moreover, the weak dependence of $\tau_m$ on excitation regime indicates that ultrafast transport properties in disordered silicon are controlled predominantly by structural disorder rather than by strong-field excitation pathway. This robustness enhances the interpretability of pump–probe experiments in realistic materials.

\section*{Outlook}

The present analysis neglects phonons and electron–electron scattering, restricting attention to sub-30-fs dynamics. On longer time scales, additional irreversible channels would further reduce coherence and modify transport. Incorporating lattice motion within time-dependent frameworks would allow exploration of the crossover between purely electronic dephasing and phonon-assisted relaxation.

Nevertheless, within the ultrafast window considered here, the combined TDDFT–Drude–SBE framework provides a consistent and quantitatively supported description of how structural disorder governs electronic dynamics in silicon.